
\chapter{יישומים ומסקנות נוספות}

\section{יישום לעיבוד \textenglish{ECG}}

יישמנו את המודל לניפוי רעש ממדידות \textenglish{ECG} (\textenglish{Electrocardiogram}). אות ה-\textenglish{ECG} מורכב מגלי \textenglish{P}, \textenglish{QRS} ו-\textenglish{T} עם תדרים ספציפיים. רעש \textenglish{EMG} (שרירים), \textenglish{baseline wander} ורעש חשמל בתדר \textenglish{50/60 Hz} מקשים על ניתוח. הדמדל מושב ל-\textenglish{ECG} הציג שיפור של \textenglish{8.4 dB} ב-\textenglish{SI-SNR} עם אימון של \textenglish{10,000} צמדי אות נקי/רועש מ-\textenglish{PhysioNet}.

\section{כיווני שיפור אדריכלי}

שלושה שיפורים מוצעים לגרסה הבאה:

\begin{enumerate}
\item \textbf{\textenglish{Conformer} במקום \textenglish{BiLSTM}}: ארכיטקטורת \textenglish{Conformer} \cite{park2024hybrid} המשלבת קונבולוציה עם \textenglish{attention} הדגימה עדיפות בחילוץ אותות קוליים. ניסויים ראשוניים שלנו הראו שיפור של \textenglish{0.6 dB} נוסף.
\item \textbf{מנגנון \textenglish{Masking} ספקטרלי}: שינוי הפלט מחיזוי פרמטרים ישירים למסיכה ספקטרלית (\textenglish{spectral mask}) מפחית שגיאות פאזה בתנאי רעש קיצוניים.
\item \textbf{פירוק לא מפוקח (\textenglish{Unsupervised Decomposition})}: שימוש ב-\textenglish{VAE} (\textenglish{Variational Autoencoder}) לדחיסת הייצוג הסמוי עשוי לאפשר הכללה טובה יותר לתדרים שלא נראו באימון.
\end{enumerate}

\section{תקציר מחקר}

מחקר זה תרם מסגרת \textenglish{BiLSTM} מאומתת לחילוץ גלי סינוס מאותות רועשים, עם שיפור של \textenglish{4 dB} על קו הבסיס ויישום מוצלח ל-\textenglish{ECG}. הקוד זמין בחינם תחת רישיון \textenglish{MIT} \cite{engel2017wavenet}.
